%% kapa: exam=1 round=29 year=2018 period=1 form=A subject=economics subject_ko=경제학원론 seq=12 number=52
%% source: Q-Net 감정평가사 기출문제 (cst00309) artlSeq=5202808 file=extracted/1008/5202808/2018년 제29회 감정평가사 1차시험 기출문제/2018년 제29회 감정평가사 1차시험 기출문제/1교시 A형.pdf
%% conversion: kordoc 4.24.51 -> tools/gen1.py (auto); review=proofread-sonnet; flags=none
\begin{question}{52}
소비자 甲이 두 재화 X, Y를 소비하고 효용함수는 $U(x,y)=xy$이다. X, Y의 가격이 각각 5원, 10원이다. 소비자 甲의 소득이 1,000원일 때, 효용극대화 소비량은? (단, $x$는 X의 소비량, $y$는 Y의 소비량이다.)
\begin{choices}
  \item $x=90,y=55$
  \item $x=100,y=50$
  \item $x=110,y=45$
  \item $x=120,y=40$
  \item $x=130,y=35$
\end{choices}
\end{question}
\answer{2}
